\documentclass[../cours_1re.tex]{subfiles}
\begin{document}
\setcounter{chapter}{21}
\chapter{Variance, covariance et erreur quadratique}
\label{ch:22}
\begin{objectif}
Démontrer les formules de variance ; comprendre la covariance et la meilleure approximation par une constante.
Prérequis : espérance et développement d'un carré.
\end{objectif}
\section{Mesurer une dispersion}
\begin{definition}
Si $\mu=\mathbb E(X)$, la variance est
$\operatorname{Var}(X)=\mathbb E((X-\mu)^2)$.
L'écart type est $\sigma(X)=\sqrt{\operatorname{Var}(X)}$.
\end{definition}
La variance est positive ou nulle puisque les carrés et les poids le sont.
Elle est nulle si et seulement si $X=\mu$ sur chaque issue de probabilité positive.
C'est l'expression « presque sûrement » dans un univers fini.
La variance a l'unité de $X$ au carré ; l'écart type a la même unité que $X$.

\begin{theoreme}[König--Huygens]
$\operatorname{Var}(X)=\mathbb E(X^2)-\mathbb E(X)^2$.
De plus $\operatorname{Var}(aX+b)=a^2\operatorname{Var}(X)$
et $\sigma(aX+b)=|a|\sigma(X)$.
\end{theoreme}
\begin{demonstration}
Développer $(X-\mu)^2=X^2-2\mu X+\mu^2$ et prendre l'espérance :
$\mathbb E(X^2)-2\mu^2+\mu^2=\mathbb E(X^2)-\mu^2$.
Comme $\mathbb E(aX+b)=a\mu+b$, la variable centrée associée est
$a(X-\mu)$. Son carré est $a^2(X-\mu)^2$, d'où les deux autres formules.
\end{demonstration}
\begin{exemple}
Si $X$ vaut $0,2,5$ avec probabilités $1/2,1/3,1/6$,
$\mu=3/2$, $\mathbb E(X^2)=11/2$ et
$\operatorname{Var}(X)=11/2-9/4=13/4$.
$\sigma=\sqrt{13}/2$. La valeur moyenne $1{,}5$ n'appartient pas au support.
\end{exemple}
Pour une série statistique finie, les mêmes définitions s'appliquent en prenant
les fréquences comme poids. On divise alors la somme des carrés centrés par
l'effectif total $n$. Le dénominateur $n-1$, rencontré en inférence statistique,
répond à un autre problème : estimer sans biais la variance d'une population.

\section{Additionner des variables}
\approfondissement
On définit
$\operatorname{Cov}(X,Y)=\mathbb E((X-\mathbb E X)(Y-\mathbb E Y))$.
Le développement et la linéarité donnent
$\operatorname{Cov}(X,Y)=\mathbb E(XY)-\mathbb E(X)\mathbb E(Y)$, puis
\[
\operatorname{Var}(X+Y)=\operatorname{Var}(X)+\operatorname{Var}(Y)
+2\operatorname{Cov}(X,Y).
\]
Plus généralement, le développement du carré d'une somme donne
$\operatorname{Var}(\sum_iX_i)=\sum_i\operatorname{Var}(X_i)
+2\sum_{i<j}\operatorname{Cov}(X_i,X_j)$.

Les variables $X,Y$ sont indépendantes lorsque, pour toutes leurs valeurs,
$P(X=x_i,Y=y_j)=P(X=x_i)P(Y=y_j)$.
\begin{demonstration}
Sous cette hypothèse, le calcul sur la loi jointe donne
\[
\mathbb E(XY)=\sum_{i,j}x_iy_jP(X=x_i)P(Y=y_j)
=\left(\sum_i x_iP(X=x_i)\right)\left(\sum_jy_jP(Y=y_j)\right).
\]
La covariance est donc nulle et les variances s'additionnent.
\end{demonstration}
La réciproque est fausse. Si $X$ est uniforme sur $\{-1,0,1\}$ et $Y=X^2$,
$\mathbb E(X)=\mathbb E(XY)=0$, donc la covariance est nulle.
Pourtant $P(X=0,Y=0)=1/3\ne(1/3)(1/3)$ : les variables sont dépendantes.

\section{Bernoulli et binomiale}
Une variable de Bernoulli $B$ satisfait $B^2=B$. Donc
$\operatorname{Var}(B)=p-p^2=p(1-p)$.
Pour $n$ Bernoulli indépendantes de même paramètre $p$, leur somme $X$
a pour espérance $np$ et variance $np(1-p)$, par les résultats précédents.
La marche $S_n=2X-n$ a donc pour espérance $n(2p-1)$ et variance
$4np(1-p)$. Pour $p=1/2$, sa position moyenne est nulle mais sa dispersion
augmente : écart type $\sqrt n$.

\section{Pourquoi la moyenne minimise l'erreur quadratique}
\begin{theoreme}
Pour tout réel $a$,
$\mathbb E((X-a)^2)=\operatorname{Var}(X)+(a-\mu)^2$.
L'unique constante minimisant cette erreur est donc $a=\mu$.
\end{theoreme}
\begin{demonstration}
Écrire $X-a=(X-\mu)+(\mu-a)$ et développer. Le terme croisé a pour
espérance $2(\mu-a)\mathbb E(X-\mu)=0$.
Il reste la variance et le carré annoncé, minimal uniquement lorsqu'il est nul.
\end{demonstration}
Cette propriété dépend du choix d'une erreur au carré. Une perte absolue
conduit aux médianes ; une perte asymétrique peut conduire à un autre choix.
Il faut donc définir le critère avant d'annoncer un « meilleur » estimateur.

\begin{finchapitre}
\section{Utilité et prolongements}
Variance et covariance mesurent la dispersion et les fluctuations communes.
L'identité de meilleure constante est la version probabiliste du théorème du
centre de gravité du chapitre 17.

En statistique et en apprentissage, on minimisera des erreurs quadratiques
sur des familles de prédicteurs. Les chapitres 27 et 28 donnent les premiers
cas résolus par orthogonalité ; le master étudiera aussi les limites de ces modèles.
\begin{retenir}
Centrer avant de mesurer la dispersion. Les variances s'additionnent
pour des variables indépendantes ; la linéarité des espérances, elle, est toujours valable.
\end{retenir}
\end{finchapitre}
\end{document}
